\documentclass{article}
\usepackage[utf8]{inputenc}
\usepackage{amsmath}
\usepackage{bm}

\begin{document}

In this study, the transient gas flow equation for the microporous region is solved using \textbf{explicit parameter lagging} combined with \textbf{semi-implicit temporal discretization}. To ensure convergence and accuracy, a \textbf{Picard iteration} scheme is employed for updating variables when necessary. The governing equation is expressed as:

\begin{equation}
    \frac{\partial}{\partial t}\left[\left(\phi-\frac{n_{ad}}{\rho^{ad}}\right) \rho^{g}+n_{ad}\right]+\nabla \cdot\left(\rho^{g} \mathbf{q}^{g}\right)+\nabla \cdot\left(\rho^{ad} \mathbf{v}^{ad}\right)=0
\end{equation}

To solve Equation~1, the storage and transport terms are treated separately. For the storage term, we define:

\begin{equation}
    S(p)=\left(\phi -\frac{n_{ad}}{\rho^{ad}}\right) \rho^{g}+n_{ad}
\end{equation}

Applying the chain rule, the temporal derivative of the storage term is given by:

\begin{equation}
    \frac{\partial S(p)}{\partial t}=\frac{d S(p)}{d p} \cdot \frac{\partial p}{\partial t}
\end{equation}

The derivative of $S(p)$ with respect to pressure $p$ is expanded as:

\begin{equation}
    \frac{d S(p)}{d p}=\left(\phi-\frac{n_{ad}}{\rho^{ad}}\right) \frac{d \rho^{g}}{d p}+\rho^{g}\left(-\frac{1}{\rho^{ad}} \frac{d n_{ad}}{d p}\right)+\frac{d n_{ad}}{d p}
\end{equation}

% NOTE FOR AUTHOR: Equation 4 uses \phi, but Equation 5 switches to \phi. Please verify whether these symbols represent the same quantity (porosity) and ensure consistency throughout the manuscript.

Rearranging Equation~4 yields:

\begin{equation}
    \frac{d S(p)}{d p}=\left(\phi-\frac{n_{ad}}{\rho^{ad}}\right) \frac{d \rho^{g}}{d p}+\left(1-\frac{\rho^{g}}{\rho^{ad}}\right) \frac{d n_{ad}}{d p}
\end{equation}

By introducing the isothermal compressibility coefficient for real gases, $c_g$, defined as:

\begin{equation}
    c_{g}=\frac{1}{\rho^{g}} \frac{d \rho^{g}}{d p}=\frac{1}{p}-\frac{1}{Z} \frac{d Z}{d p}
\end{equation}

Since the pressure change is small, the compressibility factor $Z$ can be assumed to be approximately constant, i.e., $\frac{dZ}{dp} \approx 0$. Consequently, Equation~6 simplifies to:

\begin{equation}
    c_{g}=\frac{1}{p}
\end{equation}

Using this result, the derivatives of gas density and adsorption capacity with respect to pressure are obtained as follows:

\begin{equation}
    \frac{d \rho^{g}}{d p}=\rho^{g} c_{g}=\frac{M p}{Z R T} c_{g}=\frac{M}{Z R T}
\end{equation}

\begin{equation}
    \frac{d n_{ad}}{d p}=\frac{n_{ad}^{\max} K}{(1+K p)^{2}}
\end{equation}

We can then define a \textbf{total storage capacity} $C(p) \equiv \frac{dS(p)}{dp}$. Substituting the expressions from Equations~8 and~9, we obtain:

\begin{equation}
    C_{micro}(p)=\left(\phi-\frac{n_{ad}^{\max} K p}{\rho^{ad}(1+K p)}\right) \frac{M}{Z R T}+\left(1-\frac{M p}{Z R T \rho^{ad}}\right) \frac{n_{ad}^{\max} K}{(1+K p)^{2}}
\end{equation}

Turning to the transport terms, nonlinearities arise from the compressibility of the free gas and the surface diffusion of the adsorbed gas. The Darcy flux for free gas and the surface diffusion velocity for adsorbed gas are given by:

\begin{equation}
    \mathbf{q}^{g}=-f(\mathrm{Kn}) \frac{k_{m}}{\mu} \nabla p
\end{equation}

\begin{equation}
    \mathbf{v}^{ad}=-\frac{D_{s}}{\rho^{ad}} \nabla n_{ad}
\end{equation}

Combining these into the divergence terms yields:

\begin{equation}
    \nabla \cdot\left(\rho^{g} \mathbf{q}^{g}\right)+\nabla \cdot\left(\rho^{ad} \mathbf{v}^{ad}\right)=\nabla \cdot\left(\rho^{g} f(\mathrm{Kn}) \frac{k_{m}}{\mu}+\frac{K n_{ad}^{\max} D_{s}}{(1+K p)^{2}}\right) \nabla p
\end{equation}

By defining the \textbf{total transmissibility coefficient} as:

\begin{equation}
    T_{micro}(p)=\rho^{g} f(\mathrm{Kn}) \frac{k_{m}}{\mu}+\frac{K n_{ad}^{\max} D_{s}}{(1+K p)^{2}}
\end{equation}

and substituting the total storage capacity $C_{micro}(p)$ and the total transmissibility coefficient $T_{micro}(p)$ into Equation~1, the governing equation for gas flow in the micropores reduces to:

\begin{equation}
    C_{micro}(p) \cdot \frac{\partial p}{\partial t}=\nabla \cdot\left(T_{micro}(p) \nabla p\right)
\end{equation}

Although this equation is more compact in form, the coefficients $C_{micro}(p)$ and $T_{micro}(p)$ remain highly nonlinear functions of pressure $p$. To achieve a linearized solution, these coefficients are ``frozen'' at a reference pressure $p^{b}$, taken here as the background pressure. Consequently, the final governing equation for gas flow in the micropores becomes linear in pressure:

\begin{equation}
    C_{micro}(p^{b}) \cdot \frac{\partial p}{\partial t}=\nabla \cdot\left(T_{micro}(p^{b}) \nabla p\right)
\end{equation}

For gas flow in the macropores, only the free gas phase is considered, and the governing equation takes the form:

\begin{equation}
    V_{i} \frac{\partial \rho_{i}^{g}}{\partial t}+\sum_{j=1}^{N_{i}} T_{ij}\left(p_{i}-p_{j}\right)=0
\end{equation}

Similarly, the storage capacity for the macropores is defined as:

\begin{equation}
    C_{macro}(p)=\frac{M}{Z(p) R T}
\end{equation}

Applying the same linearization at the background reference pressure $p^{b}$, Equation~17 is rewritten as:

\begin{equation}
    V_{i} C_{i,macro}(p^{b}) \frac{\partial p_{i}}{\partial t}+\sum_{j=1}^{N_{i}} T_{ij}\left(p_{i}-p_{j}\right)=0
\end{equation}

The fully implicit temporal discretization of Equation~16 yields:

\begin{equation}
    V_{i} C_{i,micro}\left(p^{b}\right) \frac{p_{i}^{t+\Delta t}-p_i^{t}}{\Delta t}+\sum_{j=1}^{N_{i}} T_{ij,micro}\left(p^{b}\right)\left(p_{i}^{t+\Delta t}-p_{j}^{t+\Delta t}\right)=0
\end{equation}

The transmissibility coefficient between adjacent micropores, $T_{ij,micro}\left(p^{b}\right)$, is computed using the harmonic mean:

\begin{equation}
    T_{ij,micro}\left(p^{b}\right)=\frac{a_{i} a_{j}}{a_{i}+a_{j}}, \quad \text{where} \quad a_{i}=\frac{A_{i} T_{i,micro}\left(p^{b}\right)}{d_{i}} \mathbf{n}_{i} \cdot \mathbf{f}_{i} \quad \text{and} \quad a_{j}=\frac{A_{j} T_{j,micro}\left(p^{b}\right)}{d_{j}} \mathbf{n}_{j} \cdot \mathbf{f}_{j}
\end{equation}

Analogously, the discretization of Equation~19 gives:

\begin{equation}
    V_{i} C_{i,macro}\left(p^{b}\right) \frac{p_{i}^{t+\Delta t}-p_i^{t}}{\Delta t}+\sum_{j=1}^{N_{i}} T_{ij,macro}\left(p^{b}\right)\left(p_{i}^{t+\Delta t}-p_{j}^{t+\Delta t}\right)=0
\end{equation}

The gas transmissibility coefficient between adjacent macropores, $T_{ij,macro}\left(p^{b}\right)$, is defined as:

\begin{equation}
    T_{ij,macro}\left(p^{b}\right)=\frac{f_{i}(\mathrm{Kn})+f_{j}(\mathrm{Kn})}{2} \cdot \frac{\pi R_{ij}^{4}}{8 \mu l_{ij}} \cdot \frac{\rho_{i}^{t}+\rho_{j}^{t}}{2}
\end{equation}

Finally, the transmissibility coefficient at the macropore--micropore interface is defined as:

\begin{equation}
    T_{ij}\left(p^{b}\right)=\frac{a_{i} a_{j}}{a_{i}+a_{j}}
\end{equation}

where the interface half-transmissibilities $a_i$ and $a_j$ are given by:

\begin{equation}
    a_{i}=\rho_{i}^{t} f_{i}(\mathrm{Kn}) \frac{\pi R_{i}^{4}}{8 \mu l_{i}} \qquad \text{and} \qquad a_{j}=\frac{A_{j} T_{j,micro}\left(p^{b}\right)}{d_{j}} \mathbf{n}_{j} \cdot \mathbf{f}_{j}
\end{equation}

% -----------------------------------------------------------------------
% POLISHING SUMMARY
% -----------------------------------------------------------------------
% The following changes were made to the original manuscript:
%
% 1. LATEX FORMATTING: Replaced all Markdown bold syntax (**text**) with
%    proper LaTeX \textbf{} commands. Added \usepackage{bm} for potential
%    bold-math needs.
%
% 2. SYMBOL NOTATION: Replaced the informal "fun(Kn)" notation with the
%    more standard mathematical function notation f(Kn) / f_i(Kn) throughout,
%    for consistency and readability.
%
% 3. MULTIPLICATION OPERATOR: Replaced the typographically incorrect
%    asterisk (*) in Equation 22 with the proper LaTeX \cdot operator.
%
% 4. SYMBOL INCONSISTENCY (FLAGGED, NOT MODIFIED): In Equation 4 the
%    porosity symbol is \phi, but in Equation 5 it switches to \phi.
%    A comment has been inserted above Equation 5 to flag this for the
%    author to verify and correct.
%
% 5. SENTENCE STRUCTURE & FLOW: Several sentences were restructured for
%    improved clarity, logical progression, and grammatical correctness.
%    Key improvements include:
%    - "First, regarding the storage term, let:" → more direct phrasing.
%    - "Since the pressure change is minimal" → "Since the pressure change
%      is small" (more standard phrasing in reservoir engineering).
%    - Added "Using this result," as a transition before Equations 8-9.
%    - Clarified the definition of C(p) using the \equiv symbol.
%    - "Frozen" placed in quotation marks to signal the technical usage.
%    - Replaced "we take the pressure from background pressure" with the
%      clearer "taken here as the background pressure."
%    - "The transmissibility coefficient between micropores" →
%      "between adjacent micropores" (precision).
%    - "Equation 19 is discretized as" → "Analogously, the discretization
%      of Equation 19 gives" (smoother transition).
%    - "Finally, the gas transmissibility coefficient between the macropores
%      and micropores" → "at the macropore–micropore interface" (precision).
%    - Renamed "interface terms" to "interface half-transmissibilities"
%      for physical clarity.
%
% 6. SPACING & LAYOUT: Replaced overloaded subscript notation (e.g.,
%    T_{i j, \text{micro}}) with cleaner T_{ij,micro} throughout.
%
% OVERALL ASSESSMENT: The original text presented a technically sound
% derivation of the linearized gas flow equations for a dual-porosity
% (micropore/macropore) pore-network model. The main issues were
% presentational: invalid LaTeX formatting for bold text, an informal
% function notation, an inconsistent porosity symbol, and occasional
% awkward phrasing. After polishing, the manuscript reads more fluently
% and adheres to standard LaTeX typesetting conventions. The one
% mathematical concern flagged (\phi vs. \phi) should be resolved
% by the author before submission.
% -----------------------------------------------------------------------

\end{document}
